\ssscp{Different patterns of contact with people {\&} environment}{02-03-02-04}
\ili{Sumbawa}, Flores, Sumba, \ili{Alor},
Timor, Tanimbar, Seram and \ili{Buru} are all islands ranging in size
from 2,000 square kilometres to over 30,000 square kilometres.
Many of them have significant
and very rugged mountains, and even micro-climates.
On \ili{Buru} the word for ‘mountain’ at first appears to reflect dialect
or language differences with \ili{Rana}
and \ili{Lisela} \it{fude} ‘mountain’ in the north,
and \ili{Masarete} \it{fuka} ‘island,
mountain’ in the south.
In fact, it is a reflection of different geology on the north and south coasts.
The island has a sharp geological southern uplift with cliffs
and tall mountains rising strait from the sea,
and \it{fuka} refers to mountains made of uplifted coral (\frf{02-23} and \frf{02-24}).
The north coast has a narrow coastal plain
and gentler hills composed of soil (see \frf{02-25});
these are \it{fude}.

\begin{figure}
  \subfigure[\it{fuka}: uplifted  coral in south \ili{Buru}]{\label{fig:02-23}
    \frame{
\includegraphics[width=0.31\textwidth]{02-23.jpg}}}\hspace{2mm}
  \subfigure[\it{fuka}: uplifted  coral in south \ili{Buru}]{  \label{fig:02-24}
    \frame{\includegraphics[width=0.31\textwidth]{02-24.jpg}}}\hspace{2mm}
  \subfigure[\it{fude}: mountains of soil in north \ili{Buru}]{\label{fig:02-25}
    \frame{\includegraphics[width=0.31\textwidth]{02-25.jpg}}}
\fcp{Hills and mountains in Buru}{02-23--02-23}
\end{figure}

In addition to synonyms created by taboos described above,
and different environmental features,
there are also dialectal “synonyms” created by different patterns of contact.
For example, \ili{Lisela} in the north of \ili{Buru} (a \tsc{Sula-Buru} language)
has had intense contact over centuries with the recently extinct
\ili{Kayeli} language spoken in the eastern part of \ili{Buru} (an \tsc{\ili{Ambon}-Seram} language).
And for several centuries the king of \ili{Lisela} was subordinate
to and paid tribute to the king of \ili{Kayeli}.
Additionally, there are around 12,000 immigrants from Sula island
to the north, settled along the north coast of \ili{Buru} for 12 or more generations
\citep{Grimes1991c,Grimes2000a,Grimes2010b}.
They have no impact on the south.

PMP *p = \it{p} in \ili{Buru}
and \ili{Lisela} native vocabulary,
as in PMP *pulut ‘birdlime,
sticky sap of breadfruit’ becoming \ili{Buru} \it{polo-n}
and \ili{Lisela} \it{pulu-n} ‘glue,
sticky substance, starchy paste’;
and PMP *sapa ‘what?’ becoming \ili{Buru}/\ili{Lisela} \it{sapa-n} ‘what?’.
But in \ili{Kayeli} PMP *p > \it{h}.
So PMP *paRih ‘skate,
stingray’ becomes \ili{Buru} \it{pahi} ‘stingray’
and \ili{Kayeli} \it{hali} ‘stingray’.
PMP *pajay ‘rice’ becomes \ili{Buru} \it{pala} ‘rice’
and \ili{Kayeli} \it{hala} ‘rice’.
PMP *qapuR ‘lime, calcium’ becomes \ili{Buru} \it{apu} ‘lime,
chalk, calcium’ and \ili{Kayeli} \it{ahul} ‘lime, chalk, calcium’.
So \ili{Lisela} \it{hali} ‘stingray’,
\it{hala} ‘rice’, and \it{ahul} ‘lime, calcium’ identify themselves
as \ili{Kayeli} loans, rather than inherited reflexes of the PMP forms.
On a mountainous island of nearly 10,000 square kilometres,
people in the south of \ili{Buru} know \it{apu}, but not \it{ahul}.
People in the north of \ili{Buru} know \it{ahul}, but not \it{apu}.
The \ili{Rana} dialect in the middle celebrates its links through marriage
both to the north and south.
So it is not surprising that the sung ballad of Raja Bau with
195 stanzas in poetic parallelisms uses both \it{apu}
and \it{ahul} as synonyms.

Like many Austronesian languages,
\ili{Buru} has phrase level genitive enclitics
(at the end of the word
to which it attaches).
So \it{(geba) faha=n} ‘(person) hand/arm-\tsc{3sg.gen}’ (< PMP
*qabaRa ‘shoulder’), \it{faha=ŋ} ‘my arm’, \it{faha=m} ‘your arm’,
\it{pani=n} ‘its wing’ (< PMP *panid ‘wing’),
\it{aa=n} ‘his jaw’ (< PMP *qazay ‘chin’), etc.
Sula, an island to the north with thousands of immigrants on
the north coast of \ili{Buru},
many of who say their line has been there for 12+ generations,
has preposed genitives, as in \it{n-aya} ‘his jaw’.
Both \ili{Buru} and Sula (\ili{Sanana}) use their genitives
in vocative constructions (direct address).
So a \ili{Buru} person (\ili{Rana} dialect) addresses his parents as \it{ama=ŋ}
‘my father’ and \it{ina=ŋ} ‘my mother’ (< PMP *ama ‘father’; *ina ‘mother’),
whereas a Sula person addresses his parents as \it{ŋ-ama} ‘my
father’ and \it{ŋ-ina} ‘my mother’.
The \ili{Lisela} vocatives are the same as the Sula forms,
indicating they are Sula loans.
The languages in the Sula Islands are the only ones with preposed
genitives with which \ili{Lisela} has had significant contact.

There are also many non-cognate forms that differ in different
dialects, some triggered by different patterns of contact,
some by tabooing, and the motivation for others remains unclear.
So ‘cuscus’ is \it{tonal} in the \ili{Masarete} dialect of \ili{Buru} (the
final \it{l{\#}} marks this word as unusual in its structure,
suggesting it is a loan from a south\ili{east} Sulawesi source -- compare
\ili{Kambowa} \it{tonali} ‘dwarf cuscus’ [\it{Strigocuscus celebensis}]),
\it{lemet} in the Wae Sama dialect,
and \it{blafen} ‘cuscus (generic)’
in the \ili{Lisela} language.
‘Mosquito’ is \it{inhadat} in the \ili{Masarete} dialect (a taboo synonym),
\it{seŋet} in the \ili{Rana}
and Wae Sama dialects (PMP retention with a shift in meaning),
and \it{mimaun} in the \ili{Lisela} language.
We find similar variation for skin,
lips, wing, megapode, thatch,
fish trap, wave, sleep, wake up, hungry,
cold/flu, defecate, good, many,
husband, wife, cloth, machete,
medicine, quick, some verbs of speaking,
some prepositions, and many other things.
Why would they need competing terms to talk about the same thing
on the same island? The simple fact is,
they do have competing terms for many things
and for a variety of reasons.

Some languages may appear to have ``competing words'' that actually
refer to different species or varieties.
For example, \ili{Lisela} (\tsc{Sula-Buru}) has \it{blafen} ‘cuscus
(generic)’ and \it{ton gawat}
‘referring to the spotted cuscus (\it{Spilocuscus maculatus})’.

Similar variation is found in other larger languages scattered
across mountains and coasts, and having different areas with different contact
patterns, such as \ili{Alune}, \ili{Wemale},
\ili{Manggarai}, \ili{Sika}, \ili{Lamaholot}, \ili{Kambera},
the \ili{Meto} languages, the Rote languages,
\ili{Tetun}, \ili{Mambae}, and \ili{Yamdena}.

So, when \citet[268]{Blust2012b} says, “we must ask why a language
would have synonyms for ‘cuscus’”,
the response is fairly simple,
“There are many reasons,
and they are all very well documented in the region.” Appealing
to the situation in Hawai‘i where the Austronesians encountered
new flora and fauna in a context of no prior human habitation,
and where seeking loan sources is futile,
\citet[273]{Blust2012b} boldly declares about the dispute over
marsupial terms in eastern Indonesia that,
“there was also no reason for \tsc{An} immigrants to borrow the
name of the cuscus rather than invent one from scratch.”

Yet the Hawai‘i situation is not at all parallel to the situation
in our target region.
\ili{As} we have shown (\srf{02-01}),
there is evidence of over 40,000 years of human habitation.
The Austronesian-speaking peoples did not enter a previously uninhabited region.
In \trf{02-Features} we gave several examples in which Austronesian
languages borrowed Papuan ways of saying things when they already
had perfectly good ways of saying them in their existing communication
framework, such as ‘hair {\tl} leaf of head’,
‘liver {\tl} insides’, ‘banana’,
the preposed possessor, alienable-inalienable distinctions,
clause final negation, and so forth.
Why invent a word,
when other people that you interact with,
who are not enemies,
already have names for it? Coining a new word “from scratch”
is usually quite rare.

